\subsection{Trade-off model 3--4B tanpa contoh}

Pada keenam kombinasi 3--4B, tidak ada konfigurasi yang sekaligus ber-TPR
tinggi dan ber-FPR rendah (Tabel~\ref{tab:baseline-models}). Kedua ujungnya
ditempati model yang berbeda: Qwen2.5-Coder-3B DIFF mencapai TPR 84,54\%
tetapi menandai hampir separuh pembaruan tipikal, sedangkan Qwen3-4B DIFF+I
menahan alarm tipikal di 3\% dengan melewatkan lebih dari sepertiga MAL
berlabel sumber. Pada protokol ini, menaikkan satu sisi menurunkan sisi lain,
dan itulah titik yang hendak diuji contoh pembanding serta profil intent pada
rancangan utama.

Blok intent sendiri belum terbukti menjadi pengungkit yang stabil. Arahnya
berbeda antar model, dan pada Phi-4-mini kenaikannya justru datang dari
transisi keluaran invalid menjadi valid, bukan dari pasangan yang kedua
responsnya valid. Selama keluaran invalid dihitung benign, validitas format
dan kemampuan menilai kode bercampur di dalam satu angka. Ablasi intent yang
konfirmatori karena itu perlu dibaca bersama validitas JSON, bukan dari
selisih TPR saja.

Seluruh pembacaan di atas memakai label sumber tanpa gold manusia, tanpa
subset \emph{hard negative}, dan tanpa interval kepercayaan, sehingga sifatnya
menunjuk arah dan bukan estimasi.

\subsection{Kecukupan konteks dan batas generalisasi}

Diagnosis dependency menunjukkan bahwa keputusan model dapat bergantung
pada kode yang tidak terlihat dalam diff target. Namun, kedua studi dependency
tidak mendukung anggapan bahwa penambahan konteks selalu memperbaiki deteksi.
Pada satu keluarga dependency, perubahan snapshot mengubah seluruh verdict;
pada kohort dependency lain, tambahan isi tidak mengubah keputusan dan
menambah penggunaan token. Temuan ini memotivasi pemeriksaan relevansi sumber,
versi yang terpasang, dan cakupan kode sebelum memperluas konteks.

Granite 3.3 8B juga tidak mengubah keputusan pada kohort yang sama. Model yang
lebih besar dari keluarga lain tidak otomatis mengatasi keterbatasan
informasi dalam masukan. Namun, hanya satu target MAL dan 14 kandidat HN
yang diuji, sehingga hasil ini tidak menunjukkan bahwa Granite dan Qwen
setara pada populasi pembaruan npm. Perluasan model harus disertai kohort
yang cukup beragam dan label yang dapat dipertanggungjawabkan.

Batas validitas utama adalah seleksi post-hoc pada studi versi dependency,
keterpaparan dev, ketergantungan antartarget yang berbagi dependency,
ketiadaan gold manusia, serta cakupan HN studi isi dependency yang belum
mencakup perilaku runtime sensitif.
Kedua studi menggunakan konteks kode dependency, bukan retrieval contoh KB
berlabel. Hasilnya tidak menjawab hipotesis keunggulan contoh benign sebagai
pembanding malicious, maupun efek intent terhadap FPR pada studi utama.

Evaluasi berikutnya perlu mencakup beberapa keluarga dependency dengan
provenance temporal, memisahkan versi resolved dari versi bersyarat, dan
menyertakan pembaruan sah dengan perilaku network, process, atau lifecycle
yang diverifikasi independen. Card KB harus ditulis dan diperiksa oleh dua
anotator; target, dependency terkait, dan klon yang telah terpapar perlu
dikecualikan dari KB dan test final. Audit ID bukti perlu dilengkapi penilaian
apakah kutipan mendukung alasan, karena format yang valid tidak menjamin
ketepatan penalaran.
